\scanpage{78}
\deleted{mod $2\omega_1,2\omega_2$ deux à deux}. $l_i$ la multiplicité de $a_i$. Alors $g(z)=f(z)\prod_{1\leq i\leq k}\bigl[\unsure{\wp|p}(z)-\wp(a_i)\bigr]^{l_i}$ est fonction elliptique paire, n'ayant plus de pôles que dans $\unsure{\mathfrak P|\mathfrak B}$. Il suffit donc d'appliquer c).

\deleted{e) Toute fonction elliptique impaire $f$ est égale à $\wp'g$, où $g$ est une fonction elliptique paire}

e) \uline{Le corps $E$ des fonctions elliptiques est de degré 2 sur $C\langle\wp\rangle$}. En effet, $C\langle\wp\rangle$ est le corps des invariants de l'automorphisme $f\to\unsure{\check f|\breve f}$ de $E$, qui est d'ordre 2 (noter que $E\ne C\langle\wp\rangle$, car $\wp'$ est \uline{impaire})

comme corollaire, puisque $\wp'\in\complement C\langle\wp\rangle$,

f) \uline{Toute fonction elliptique $f$ se met d'une façon et d'une seule sous la forme}
\[
f=R_1(\wp)+\wp'R_2(\wp)\tag*{(3)}
\]
\uline{où $R_1$ et $R_2$ sont des fractions rationnelles}

\deleted{\unsure{Énonçons|d'ailleurs}} \inserted{Notons encore}

g) \deleted{\unsure{La totalité des pôles|Le nombre de pôles}} \uline{Pour que la fonction $f$ dans (3) ait ses pôles contenus dans $\mathfrak P$, il faut et il suffit que $R_1$ et $R_2$ soient des polynômes ; alors $0$ est au moins pôle double}.

\deleted{Il suffit de prouver que si les pôles de $f$ sont tous dans $\mathfrak P$, $f$ est de la forme $P_1(\wp)+\wp'\unsure{R_2|P_2}(\wp)$, $P_1$ et $P_2$ polynômes. Ecrivant $f=\frac{f+\check f}{2}+\frac{f-\check f}{2}$, on est ramené au cas où $f$ est paire -- traité dans c) -- resp\textsuperscript{t} où $f$ est impaire. Mais \deleted{sur} (3) on voit que cela signifie que $f$ est}

% archive sheet number: 77
